\section{Transporting the retained counterexample mechanism}
\label{sec:retained-mechanism}

This section starts from the explicit polynomial used in the accompanying
fluid reconstruction, not from a search over unrelated test functions.
Tao's displayed map and geometric explanation\cite{tao2026jacobian},
Speyer's marked-factor discussion\cite{speyer2026jacobian}, and
Poplett's complete polynomial-coordinate proof\cite{poplett2026alias}
identify the source. The calculations below retain its original coefficients.
The interfaces developed here are the full inverse-fibre cubic, its coefficient
heat flow, and the conductor model obtained from its two-sheet fibre product.
These calculations establish precise local and dynamical connections;
they do not identify a zero of the Riemann zeta function off its critical line.

\subsection{Original coordinates and the entire inverse fibre}
Use spatial coordinates $(x,y,w)$: $w$ is the third source coordinate
called $t$ in the cited polynomial, with the explicit identity renaming
$(x,y,t)\mapsto(x,y,w)$. Evolution time below is $\tau$ or $s$.
Write the target as $(a,b,c)$ and set
\begin{align}
 F_1&=(1+xy)^3w+y^2(1+xy)(4+3xy),\notag\\
 F_2&=y+3x(1+xy)^2w+3xy^2(4+3xy),\notag\\
 F_3&=2x-3x^2y-x^3w. \label{eq:retained-F}
\end{align}
Here is a short exact determinant check. On $x\ne0$, put
$v=1+xy$, $h=2-3xy-x^2w$, $A=v^2+v-v^3h$ and
$B=4v+2-3v^2h$. Then $F=(A/x^2,B/x,xh)$ and
\[
 \det\frac{\partial(x,v,h)}{\partial(x,y,w)}=-x^3,
 \qquad
 \det\frac{\partial F}{\partial(x,v,h)}
 =x^{-3}\det\begin{pmatrix}
 -2A&2v+1-3v^2h&-v^3\\
 -B&4-6vh&-3v^2\\h&0&1
 \end{pmatrix}=2x^{-3}.
\]
The numerator in the last equality expands to
$(4-6vh)(-2v^2-2v+3v^3h)
+(2v+1-3v^2h)(4v+2-6v^2h)=2$.
The product is $-2$. Since $\det DF$ is a polynomial, the identity on
$x\ne0$ proves $\det DF=-2$ everywhere, including $x=0$.

\begin{proposition}[An exact inverse-fibre chart]
\label{prop:retained-fibre}
For fixed $(a,b,c)\in\C^3$, define
\begin{equation}
 P_{a,b,c}(r)=cr^3-2r^2+br-2a,
 \qquad \alpha(r)=\frac12P'_{a,b,c}(r).
 \label{eq:retained-cubic}
\end{equation}
The points in $F^{-1}(a,b,c)$ with $x\ne0$ correspond bijectively
to the simple roots $r$ of $P_{a,b,c}$, by
\begin{equation}
 x=\alpha^{-1},\qquad y=r-\alpha,\qquad
 w=5\alpha^2-3r\alpha-c\alpha^3.
 \label{eq:retained-inverse}
\end{equation}
When $c=0$ there is one additional point
$(0,b,a-4b^2)$. When $c\ne0$ there is no point with $x=0$.
\end{proposition}
\begin{proof}
At a point with $x\ne0$ define $r=y+1/x$ and $\alpha=1/x$.
Solving $F_3=c$ gives the formula for $w$ in
\eqref{eq:retained-inverse}. Substitution in the other two original
polynomials gives, before imposing their target values,
\[
 F_2=4r+2\alpha-3cr^2,\qquad
 F_1=r^2+r\alpha-cr^3.
\]
Thus $F_2=b$ is precisely $2\alpha=3cr^2-4r+b=P'(r)$.
After this substitution, $2(F_1-a)=P(r)$. These identities prove
both directions, and recovering $r=y+1/x$ proves injectivity of the
root parametrization. The requirement $\alpha\ne0$ is exactly
the simple-root requirement; a multiple root is not a finite point
of this chart. Finally $F(0,y,w)=(w+4y^2,y,0)$, which proves the
separate boundary statement without division by $x$.
\end{proof}

For $c\ne0$ the cubic discriminant is
\begin{equation}
 \operatorname{Disc}_r P
 =4\bigl(b^2-cb^3+18abc-16a-27a^2c^2\bigr).
 \label{eq:retained-discriminant}
\end{equation}
This follows by substitution of $(c,-2,b,-2a)$ into the expanded
cubic discriminant, or by eliminating $r$ between $P$ and $P'$.
The same polynomial remains defined at $c=0$, but the polynomial
then has degree two and the extra $x=0$ point must still be included.
For example $(a,b,c)=(-1/4,0,0)$ gives $r=\pm1/2$ and
$\alpha=-2r$, yielding $(1,-3/2,13/2)$ and
$(-1,3/2,13/2)$; the additional point is $(0,0,-1/4)$.
At $(a,b,c)=(0,0,0)$ the quadratic has a double root, hence neither
of these two chart points survives; the additional point is $(0,0,0)$.
In particular, loss of the two inverse branches is not a zero of
$\det DF$ at a finite source point.

\subsection{The same branch as an incompressible trajectory}
Let $J=DF$ on $\R^3$ and define
\begin{equation}
 W=J^{-1}e_1=-\tfrac12\nabla F_2\times\nabla F_3
   =\operatorname{curl}(-\tfrac12F_2\nabla F_3).
 \label{eq:retained-piola}
\end{equation}
The scalar triple products with $\nabla F_i$ prove $JW=e_1$;
the determinant computation shows that its coefficients are
polynomial. Commuting mixed derivatives gives
$\operatorname{div}W=0$.
For $s<1/4$ define
\[
 \gamma_\pm(s)=\left(\pm(1-4s)^{-1/2},
       \mp\tfrac32(1-4s)^{1/2},\tfrac{13}2(1-4s)\right).
\]
Both satisfy $F(\gamma_\pm(s))=(s-1/4,0,0)$ by the inverse chart.
Differentiation gives $J\gamma_\pm'=e_1$, hence
$\gamma_\pm'=W(\gamma_\pm)$. Their first coordinate escapes at
$s=1/4$. The inverse-root coordinate on the positive branch is
$r=-\sqrt{1-4s}/2$, with $r^2=1/4-s$ and $r'=-1/(2r)$.

The metric making this transport an isometry is explicitly
$g=J^{\mathsf T}J$, with $\det g=4$.
Locally $X=F(x,y,w)$ is a coordinate system and $g=\sum dX_i^2$.
Thus $W$ is the pullback of the constant field $e_1$, is parallel,
and has $g$-norm one. It solves the intrinsic incompressible viscous
equation for every positive viscosity with constant pressure and the
Hodge Laplacian on one-forms. Its displayed trajectories have finite
$g$-length before escaping, proving metric incompleteness: a sequence
of their times tending to $1/4$ is metric Cauchy but has no finite
coordinate limit. A smooth positive metric gives the ordinary local
topology, as follows by bounding its eigenvalues above and below on
a small closed coordinate ball.

This transport retains its metric, rather than silently replacing it
by the original Euclidean one. For a scalar target function $f$, set
$W_i=J^{-1}e_i$. The chain rule proves
\begin{equation}
 W_i(f\circ F)=(\partial_i f)\circ F,
 \qquad \sum_iW_i^2(f\circ F)=(\Delta_Xf)\circ F.
 \label{eq:retained-operators}
\end{equation}
The fields commute: their bracket annihilates each $F_j$, and
invertibility of $J$ then makes that bracket zero.
The ordinary source Laplacian is not the operator in
\eqref{eq:retained-operators}: $\Delta_{x,y,w}F_1(0)=8$ whereas
$\Delta_X X_1=0$. These are explicit operator coordinates for the
fluid transfer, not a conclusion that the constructions are unrelated.

\subsection{The released fluid amplifier in material coordinates}
The released Boussinesq construction\cite[Section~3.1, Lemma~3.1]{ab2026boussinesq}
provides the following local wave. We retain its profile rather than
replace it by a different trial wave. Write $J_2(a,b)=(-b,a)$,
$\kappa=|\zeta|^2$, $\zeta(t_0)\ne0$, $\lambda>0$, and
$s=\lambda\zeta\cdot x$.
The earlier fields are $u_{\rm old}=Dx$, $\theta_{\rm old}=G\cdot x$,
with $\operatorname{tr}D=0$. A smooth odd periodic profile $f$ has
an even periodic primitive $p$, $p'=f$. The new fields are
\[
 \vartheta=\Theta f(s),\quad
 \psi=\frac{\Omega}{\lambda^2\kappa}p(s),\quad
 v=\frac{\Omega}{\lambda\kappa}J_2\zeta f(s),\quad
 \varpi=\Omega f'(s).
\]
Here $v=J_2\nabla\psi$, $\varpi=\operatorname{curl}v=\Delta\psi$.
The equations for the inviscid scalar and vorticity have right-hand
sides $0$ and $\partial_1\theta$, respectively, before the specified
forcing residuals. Their exact increments vanish when
\begin{equation}
 \dot\zeta=-D^{\mathsf T}\zeta,\qquad
 \dot\Theta=-\frac{J_2\zeta\cdot G}{\lambda\kappa}\Omega,
 \qquad \dot\Omega=\lambda\zeta_1\Theta.
 \label{eq:retained-fluid-amplifier}
\end{equation}
For completeness, $\dot\zeta=-D^{\mathsf T}\zeta$ makes
$(\partial_t+Dx\cdot\nabla)s=0$. Since $J_2\zeta\cdot\zeta=0$,
$v\cdot\nabla\vartheta=v\cdot\nabla\varpi=0$.
The old vorticity is spatially constant. The scalar residual is
$[\dot\Theta+\Omega(J_2\zeta\cdot G)/(\lambda\kappa)]f(s)$;
the vorticity residual is $[\dot\Omega-\lambda\zeta_1\Theta]f'(s)$.
These expressions prove every cancellation in
\eqref{eq:retained-fluid-amplifier}. They use no sinusoidal identity.

The covector map has an exact inverse-flow origin. If
$\dot M=DM$, $M(t_0)=I$, then
$\zeta=M^{-\mathsf T}\zeta(t_0)$, as differentiation proves.
Jacobi's determinant identity gives
$\partial_t\det M=(\operatorname{tr}D)\det M=0$, so $\det M=1$.
More generally a smooth incompressible material flow
$\dot X(a,t)=u(X(a,t),t)$ with $X(a,t_0)=a$ has
\[
 \partial_t(D_aX)=(D_xu)\circ X\,D_aX,
 \quad \det D_aX=1,
 \quad \nabla_x(\phi_0\circ X^{-1})\circ X
       =(D_aX)^{-\mathsf T}\nabla_a\phi_0.
\]
The first identity follows by differentiating the flow in its label;
the other two follow by the determinant identity and chain rule.
These formulas exhibit exactly how gradient information is transported
while the determinant stays fixed. The particular amplification and
reset sequence, compact localization, and all-derivative forcing control
belong to the source construction, not to the determinant identity alone.

One can retain viscosity directly in this calculation. Scalar
diffusivity $\eta\ge0$ and vorticity viscosity $\nu\ge0$ add
$\eta\lambda^2\kappa\Theta f''$ and
$\nu\lambda^2\kappa\Omega f'''$ to the two right-hand sides.
For the special profile $f(s)=\sin s$ and stationary background
$D=0$, $G=-Ae_2$, $A>0$, put
$\zeta=r(\sin\phi,\cos\phi)$ with $r>0$, and $k=\lambda r$.
Then the amplitude matrix and its two eigenvalues are exactly
\begin{align*}
 B_{\eta,\nu}&=\begin{pmatrix}
 -\eta k^2&A\sin\phi/k\\k\sin\phi&-\nu k^2
 \end{pmatrix},\\
 \mu_\pm&=-\frac{(\eta+\nu)k^2}{2}
 \mathbin{\pm}\sqrt{\frac{(\eta-\nu)^2k^4}{4}+A\sin^2\phi}.
\end{align*}
Expanding $\det(\mu I-B)$ proves the formula; comparing the
nonnegative square root with $(\eta+\nu)k^2/2$ gives
$\mu_+>0$ exactly when $A\sin^2\phi>\eta\nu k^4$.
For the actual localized profile, the displayed $f''$ and $f'''$
terms remain; there is no replacement of those functions by $-f,-f'$.

There is likewise an exact, reversible force map in the original
three-dimensional Cartesian coordinates. From an Euler triple
$\partial_tu+(u\cdot\nabla)u+\nabla p_E=f_E$, $\operatorname{div}u=0$,
all Navier--Stokes triples with the same velocity and viscosity $\nu$
are
\[
 p_{NS}=p_E+q,\qquad f_{NS}=f_E-\nu\Delta u+\nabla q,
 \qquad \operatorname{curl}f_{NS}
   =\operatorname{curl}f_E-\nu\Delta\operatorname{curl}u.
\]
Substitution proves sufficiency; subtraction proves necessity.
The curl identity follows from commuting distributional derivatives,
so it retains every pressure adjustment. The companion Euler
manuscript\cite[Theorem~1.1 and Sections~3--4]{ab2026euler}
is the source of the three-dimensional material-layer construction.
This section reconstructs the amplifier and its exact transfer terms;
it does not certify that manuscript's full multiscale proof or replace
the viscous transfer by dropping $\nu\Delta u$.

\subsection{Exact heat transport of the fibre cubic}
Keep $a_0,b_0,c$ fixed and move the target by
\begin{equation}
 a(\tau)=a_0+2\tau,\qquad b(\tau)=b_0+6c\tau.
 \label{eq:retained-coefficient-flow}
\end{equation}
For $P(\tau,r)=P_{a(\tau),b(\tau),c}(r)$ direct differentiation gives
\begin{equation}
 P_\tau=6cr-4=P_{rr}.
 \label{eq:retained-heat}
\end{equation}
Thus the original inverse-fibre family is closed under the forward
heat operator, with every coefficient and the time sign specified.
On a simple root, $r_\tau=-P_{rr}/P_r$. The velocity of its lifted
point in the original source coordinates is
$J^{-1}(2,6c,0)^{\mathsf T}$, by differentiating its target.
No such finite lift is assigned when $P_r=0$.

With Newman time $t=-\tau$, write
$\mathcal P(t,r)=P_{a_0-2t,b_0-6ct,c}(r)$. Then
$\mathcal P_t=-\mathcal P_{rr}$ and
$r_t=\mathcal P_{rr}/\mathcal P_r$, the same local zero equation
as in Section~\ref{sec:rh-newman-collision}.
For $b_0=c=0$ this becomes exactly
\[
 \mathcal P(t,r)=-2\bigl(r^2-2(t-t_c)\bigr),\qquad t_c=a_0/2.
\]
For real $a_0$ and real $t$, the double root occurs at $(t_c,0)$;
the two roots are real for
$t>t_c$ and nonreal for $t<t_c$. The point $a_0=-1/4$ of the
original three-point collision corresponds to $t_c=-1/8$ in this
particular heat orbit. Neither shifting this model time origin nor
choosing another $a_0$ changes the distinguished initial function
$H_0(z)=\xi(1/2+iz/2)/8$ of the actual Newman family.

There is also an exact nonlinear velocity equation. Wherever $P\ne0$,
set $u=-2P_r/P$. Writing locally $L=\log P$ gives
$L_\tau=L_{rr}+L_r^2$. Differentiating this equality proves
$u_\tau+uu_r=u_{rr}$. A zero of multiplicity $m$ gives
$u=-2m/(r-r_*)+O(1)$ by factoring $P=(r-r_*)^mA$, $A(r_*)\ne0$.
This is the singularity of a specified logarithmic-derivative map.

The arithmetic transport can already be written without guessing a
new zero. Define
\[
 \widehat H(t,x,y,w)=H_t(F_1(x,y,w)).
\]
The established Newman equation and
\eqref{eq:retained-operators} give
\begin{equation}
 \partial_t\widehat H=-W_1^2\widehat H.
 \label{eq:retained-arithmetic-pullback}
\end{equation}
Its values are constant on every fibre of $F$. In particular the
nonzero functional
$D_s=\delta_{\gamma_+(s)}-\delta_{\gamma_-(s)}$ annihilates every
scalar $f\circ F$, including $\widehat H(t,\cdot)$, while
$D_s(x)=2(1-4s)^{-1/2}$ diverges. This proves exactly which branch
information scalar descent loses and supplies the retained
branch-sensitive observable. The cancellation is not a negative
Weil value or an off-critical zero.

Even a local replacement of $H_t$ by a polynomial factor must carry
its analytic unit. The exact product identity is
\begin{equation}
 (\partial_t+\partial_z^2)(AP)
 =A(\partial_t+\partial_z^2)P+(A_t+A_{zz})P+2A_zP_z.
 \label{eq:retained-unit-defect}
\end{equation}
At a simple zero of $P$, the ratio $(AP)_{zz}/(AP)_z$ is
$P_{zz}/P_z+2A_z/A$. Thus discarding the unit would change the
zero dynamics by an explicitly known term. The present reconstruction
transfers the mechanism and its equations; it does not replace the
arithmetic initial datum by a chosen cubic.

\subsection{Two inverse sheets produce a conductor crossing}
\label{sec:retained-conductor}
On $b=c=0$ the inverse-root cover is the map
$\pi:\C_r\to\C_a$, $a=-r^2$. Its two-fold self-fibre product has
coordinate ring
\[
 B=\C[r_1,r_2]/(r_1^2-r_2^2).
\]
The linear changes
\begin{equation}
 u=r_1-r_2,\quad v=r_1+r_2,
 \qquad r_1=(u+v)/2,\quad r_2=(v-u)/2
 \label{eq:retained-node-map}
\end{equation}
give the exact ring isomorphism $B\simeq R=\C[u,v]/(uv)$.
The component $u=0$ records equal roots, and $v=0$ records opposite
roots. For the physical $x\ne0$ inverse chart one must remove $r_1r_2=0$;
the crossing at the origin belongs to its algebraic completion.
Adjoining a free coordinate $z$ gives the corresponding local surface
$uv=0$ in three-space. This is an actual map from the retained
polynomial to the double-crossing model used in the analysis of
Campana--Demailly--Peternell\cite[Section~7, p.~692]{cdp2020},
through a stated fibre-product construction.

We first compute in the curve ring $R$; the free surface coordinate
will be restored below. The normalization ring is
$\widetilde R=\C[u]\oplus\C[v]$.
The injection $R\to\widetilde R$ sends a polynomial to its
restrictions $(f(u,0),f(0,v))$. Its image consists exactly of pairs
having equal constant terms: the inverse reconstruction is
$p(u)+q(v)-p(0)$ for a pair with $p(0)=q(0)$.
The conductor is
\[
 \mathfrak c=(u,v)\subset R,
 \qquad \mathfrak c\widetilde R=u\C[u]\oplus v\C[v].
\]
Indeed a pair with zero constant terms can be multiplied by every
pair in $\widetilde R$ and remains in the image. Conversely
multiplication by $(1,0)$ forces the common constant term of any
conductor element to be zero.

The intrinsic differential module is explicitly
\begin{equation}
 \Omega_R^1=(R\,du\oplus R\,dv)/R(v\,du+u\,dv).
 \label{eq:retained-node-differentials}
\end{equation}
Its class $\theta=v\,du=-u\,dv$ is nonzero. If it were zero,
$(v,0)=h(v,u)$ in $R^2$, so $hv=v$ and $hu=0$.
The latter forces $h\in(v)$ by the just-proved normalization
description. Reducing the former in $\C[v]$ then forces $h(0,v)=1$,
contradicting $h\in(v)$. Also $u\theta=v\theta=0$.
Since $u+v$ is a non-zero-divisor (its two normalized restrictions
are nonzero polynomials), this is a genuine torsion class.
Its image in $\Omega_{\widetilde R}^1$ is zero on both branches.

In fact the kernel of $\Omega_R^1\to\Omega_{\widetilde R}^1$
is exactly $\C\theta$. For a class $A\,du+B\,dv$ to map to zero,
$A(u,0)=B(0,v)=0$, so $A=va(v)$ and $B=ub(u)$.
The relations $v^2du=u^2dv=0$ reduce this class to
$(a(0)-b(0))\theta$. This proves the kernel statement, not merely
the existence of a forgotten class. In the original root coordinates,
$\theta=r_2dr_1-r_1dr_2$: expand $vdu$ using
\eqref{eq:retained-node-map} and use $r_1dr_1=r_2dr_2$.
For the surface ring $R[z]$, differentials split as
$(\Omega_R^1\otimes_R R[z])\oplus R[z]dz$.
The last summand injects into the normalized ring's $dz$ summand,
and the preceding calculation works coefficient by coefficient in $z$.
Thus its kernel is exactly $\C[z]\theta$.

The ambient restricted module is another explicit part of the same
diagram: for $S=(uv=0)\subset\C^2$ it is
$\Omega_{\C^2}^1|_S=R\,du\oplus R\,dv$.
Its nonzero conormal $d(uv)=vdu+udv$ maps to zero in
\eqref{eq:retained-node-differentials}; $vdu$ instead maps to the
nonzero class $\theta$ and subsequently to zero on normalization.
These two calculated kernels are the precise objects to retain when
comparing an ambient restricted differential with its torsion-free
quotient. They are not interchangeable vanishings. No global line
bundle, compact threefold, or CDP20 counterexample is inferred just
from this local ring: the existing global construction must transport
these classes through its own gluing maps.

Here is the global sheaf diagram in which the local calculation belongs.
Let $S$ be a reduced effective Cartier divisor in a smooth complex
manifold $X$. Put $E=\Omega_X^1|_S$, $N=\mathcal O_S(-S)$,
$T=\ker(\Omega_S^1\to j_*\Omega_{S_{\rm reg}}^1)$, and
$Q=\Omega_S^1/T$, where $j:S_{\rm reg}\hookrightarrow S$.
Let $K$ be the inverse image of $T$ under $E\to\Omega_S^1$.
The precise sequences are
\begin{equation}
 0\longrightarrow N\longrightarrow K\longrightarrow T\longrightarrow0,
 \qquad
 0\longrightarrow K\longrightarrow E\longrightarrow Q\longrightarrow0.
 \label{eq:retained-cdp-filtration}
\end{equation}
Indeed the Cartier conormal map sends a local frame to $dg$ for
a reduced defining equation $g$. At the generic smooth points of
each component, $dg$ is nonzero. A coefficient annihilating it
therefore vanishes on every component and is zero since $S$ is reduced.
This proves left injectivity in the conormal sequence. Taking the
inverse image of $T$ proves the first sequence, and its quotient proves
the second. On $uv=0$, $N$ is generated by $vdu+udv$ and $T$ by
the class calculated above.

For any line bundle $L$ on $S$, twisting these sequences is exact
because a line bundle is locally free. In particular the relevant
segments in cohomology are
\begin{align*}
0&\longrightarrow H^0(N\otimes L)\longrightarrow H^0(K\otimes L)
 \longrightarrow H^0(T\otimes L)
 \xrightarrow{\delta_L} H^1(N\otimes L),\\
0&\longrightarrow H^0(K\otimes L)\longrightarrow H^0(E\otimes L)
 \longrightarrow H^0(Q\otimes L).
\end{align*}
The connecting map is explicit: lift a torsion section to sections
$k_i$ of $K\otimes L$ on a sufficiently fine trivializing cover
where these local lifts exist; the differences $k_j-k_i$ lie in
$N\otimes L$ and form its representing cocycle.
Changing lifts adds a coboundary. In the vanishing reduction used in
the cited Section~7, setting the two endpoint groups
$H^0(N\otimes L)$ and $H^0(Q\otimes L)$ to zero leaves precisely
$\ker\delta_L$ as the ambient group $H^0(E\otimes L)$.
It does not erase $T$. This calculation identifies the exact additional
term and the actual gluing map, without assigning any global conclusion
to the local model alone.

\subsection{Retaining the two sheets in the arithmetic spectral quotient}
Connes--Consani--Moscovici\cite[Section~3.6, Proposition~3.6]{ccmProlate2024}
use the Hadamard topological ring $\mathcal H_{\le1}$ of entire
functions of order at most one. Their quotient is
\[
 Q_{\rm ar}=\mathcal H_{\le1}/
 \overline{\operatorname{span}\{\mathcal F_\mu\mathcal E(\psi_\ell^\pm)\}},
 \qquad
 (\mathcal F_\mu\mathcal E\psi_\ell^\pm)(s)=R_\ell^\pm(s)\Xi(s).
\]
Multiplication by $s$ induces an operator $S$ on this quotient;
the cited proposition identifies its spectrum with the parameters
$s$ for which $1/2+is$ is a nontrivial zeta zero. We use this
specified quotient, not an orthogonal complement and not an
unjustified replacement of the closed subspace by an algebraic ideal.
The exact relation to our heat coordinate is
\[
 \Xi(s)=\xi(1/2+is)=8H_0(2s),\qquad a=2s.
\]
Consequently the retained inverse-root cover $a=-r^2$ gives
$s=-r^2/2$, including the factor two.

Form the finite free scalar extension
\begin{equation}
 \widehat Q_{\rm ar}=\C[r]\otimes_{\C[s]}Q_{\rm ar},
 \qquad s\longmapsto-r^2/2.
 \label{eq:retained-arithmetic-extension}
\end{equation}
Every polynomial in $r$ has a unique even--odd decomposition
$p_0(-r^2/2)+r p_1(-r^2/2)$. Existence follows by separating its
even and odd powers, and uniqueness by comparing those powers.
Thus $\C[r]$ is free of rank two over $\C[s]$, and
\eqref{eq:retained-arithmetic-extension} identifies exactly with
$Q_{\rm ar}\oplus rQ_{\rm ar}$. Give it this finite direct-sum
topology. We are not asserting that this tensor product is a space
of all entire functions pulled back along $s=-r^2/2$; that would
alter their growth class.

In the ordered basis $(1,r)$, multiplication by $r$ and by the
original spectral coordinate $s$ are, respectively,
\begin{equation}
 \mathcal R=\begin{pmatrix}0&-2S\\I&0\end{pmatrix},
 \qquad \widehat S=-\tfrac12\mathcal R^2
          =\begin{pmatrix}S&0\\0&S\end{pmatrix}.
 \label{eq:retained-arithmetic-matrices}
\end{equation}
Indeed $r(q_0+r q_1)=-2S q_1+r q_0$, proving both matrices.
The sheet involution sends $(q_0,q_1)$ to $(q_0,-q_1)$;
its trace is $2q_0$ and its kernel is precisely $rQ_{\rm ar}$.
This trace is $\C[s]$-linear, not $\C[r]$-linear:
$\mathcal R(rq)=-2Sq$ lies in the even summand, so the odd
summand is not in general preserved by $\mathcal R$ itself.
The normalized trace $\pi(q_0,q_1)=q_0$ has section
$\iota(q)=(q,0)$. Therefore the sequence
\[
 0\longrightarrow rQ_{\rm ar}\longrightarrow\widehat Q_{\rm ar}
 \xrightarrow{\pi}Q_{\rm ar}\longrightarrow0
\]
is split, and $\widehat S$ preserves both summands. These are the
explicit branch-retention, branch-forgetting, and inverse maps.

The spectrum in the original arithmetic coordinate is unchanged:
$\operatorname{Spec}\widehat S=\operatorname{Spec}S$.
Here spectrum means failure of a continuous two-sided inverse in
the source quotient topology. If $S-\lambda$ has inverse $T$,
$\operatorname{diag}(T,T)$ inverts $\widehat S-\lambda$.
Conversely, for an inverse $B$ of $\widehat S-\lambda$,
$\pi B\iota$ is an inverse of $S-\lambda$, since
$\pi(\widehat S-\lambda)=(S-\lambda)\pi$ and
$(\widehat S-\lambda)\iota=\iota(S-\lambda)$.
Continuity follows from the finite direct-sum topology in both
directions. This proof carries the actual operator and its lost
summand; it does not turn a complex value of $r$ into an off-real
value of the original spectral coordinate $s$.
For an eigenvector $Sq=\lambda q$ and a scalar $\rho$ with
$\rho^2=-2\lambda$, $(\rho q,q)$ is directly an eigenvector of
$\mathcal R$ with eigenvalue $\rho$.
Thus even that change of spectral coordinate is completely explicit.

\subsection{The covariant branch derivative and its exact obstruction}
\label{sec:retained-branch-connection}
The finite extension retains two coordinates, so differentiating a section
requires retaining the derivative of the branch itself. On the punctured
base $s\ne0$, write a section as
\[
 q(s,r)=q_0(s)+r q_1(s),\qquad r^2=-2s.
\]
Since $2r\,dr/ds=-2$, one has
\[
 \frac{dr}{ds}=-\frac1r=\frac{r}{2s}.
\]
In the ordered basis $(q_0,q_1)$ define
\begin{equation}
 D=\frac{d}{ds}+A(s),\qquad
 A(s)=\begin{pmatrix}0&0\\0&\frac1{2s}\end{pmatrix},
 \qquad
 R(s)=\begin{pmatrix}0&-2s\\1&0\end{pmatrix}.
 \label{eq:retained-branch-covariant-derivative}
\end{equation}
Here $R$ is multiplication by $r$, not a new identification of the two
sheets. Direct differentiation gives the exact commutator
\begin{equation}
 [D,R]=R'+[A,R]=-R^{-1},\qquad
 R^{-1}=\begin{pmatrix}0&1\\-\frac1{2s}&0\end{pmatrix}.
 \label{eq:retained-branch-connection-commutator}
\end{equation}
For completeness, squaring the connection gives
\begin{equation}
 D^2=\frac{d^2}{ds^2}+2A\frac{d}{ds}+A'+A^2
 =\begin{pmatrix}
 \frac{d^2}{ds^2}&0\\[2pt]
 0&\frac{d^2}{ds^2}+\frac1s\frac{d}{ds}-\frac1{4s^2}
 \end{pmatrix}.
 \label{eq:retained-branch-connection-square}
\end{equation}
The inverse matrix exists precisely on the punctured base because
$R^2=-2sI$. Thus the inverse-square term is forced by the ramified map
$s=-r^2/2$, rather than being a coordinate convention.

\begin{proposition}[Exact heat transport on the two-sheet module]
Let $H_t(a)$ be the scalar heat-coordinate function already used above and
put $Z_t(s)=8H_t(2s)$. If $\partial_tH_t(a)=-\partial_a^2H_t(a)$, then
\begin{equation}
 \partial_tZ_t(s)=-\frac14\partial_s^2Z_t(s).
 \label{eq:retained-branch-heat-coordinate}
\end{equation}
On a two-component section, the corresponding branch-covariant second
derivative is $D^2$, not the componentwise operator $\partial_s^2$.
Consequently the odd component acquires the exact inverse-square term in
\eqref{eq:retained-branch-connection-square}.
\end{proposition}

\begin{proof}
The change of variable $a=2s$ gives
$\partial_a^2H_t(2s)=\frac14\partial_s^2H_t(2s)$, proving
\eqref{eq:retained-branch-heat-coordinate} after multiplication by $8$.
For a section $q_0+r q_1$, the displayed formula for $dr/ds$ gives
$dq/ds=(q_0',r(q_1'+q_1/(2s)))$, exactly $D(q_0,q_1)$.
Applying this once more gives $D^2$ and its lower-right entry in
\eqref{eq:retained-branch-connection-square}. No differentiation has
been asserted to descend to the closed quotient $Q_{\rm ar}:$ the
calculation is on the explicit punctured finite extension, with its domain
stated above.
\end{proof}

The commutator records the precise information lost by treating the branch
variable as a scalar coordinate. In particular, $D$ does not commute with
multiplication by $r$; the failure is the explicit map $-R^{-1}$. This is
the typed obstruction for transporting the material generator through the
arithmetic module. It is not a negative Weil pairing, an off-critical zero,
or a classical Navier--Stokes disproof.

For a scalar analytic unit $A(s,t)$ and a two-component section $v$, the same
connection gives the exact product rule
\begin{equation}
 \widehat L(Av)=A\widehat L v
 +\left(A_t+\frac14A_{ss}\right)v
 +\frac12 A_sDv,
 \qquad \widehat L=\partial_t+\frac14D^2.
 \label{eq:retained-branch-unit-defect}
\end{equation}
This is a coordinate identity on the punctured extension.  In particular,
if $v$ solves the homogeneous branch heat equation, the two displayed
derivative terms are the complete analytic-unit defect; they cannot be
discarded by passing to the scalar quotient.  On a simple scalar zero
$Z_t=A P$ the corresponding scalar specialization is
\begin{equation}
 \lambda'(t)=\frac14\frac{P_{ss}}{P_s}(t,\lambda(t))
             +\frac12\frac{A_s}{A}(t,\lambda(t)),
 \label{eq:retained-branch-zero-motion}
\end{equation}
on the open locus $P_s\ne0$.  At a multiple zero the quotient is undefined,
and \eqref{eq:retained-branch-unit-defect}, rather than a scalarized formula,
is the retained statement.

The scalar transported arithmetic generator from
\eqref{eq:cq-B} has a direct two-sheet lift on this punctured extension:
\begin{equation}
 \widehat B_t=\widehat S-\frac t2D,
 \qquad \widehat S=-\frac12R^2,
 \qquad [\widehat B_t,R]=\frac t2R^{-1}.
 \label{eq:retained-branch-generator}
\end{equation}
Indeed, $\widehat S$ is a polynomial in $R^2$ and hence commutes with $R$,
while \eqref{eq:retained-branch-connection-commutator} gives the last
identity. This is the exact branch-coupling term that must be retained when
the material generator is transported through the cover. It is defined on
the punctured two-sheet extension; it is not asserted to induce an operator
on the frozen quotient $Q_{\rm ar}$.

There is, however, an exact typed interface with the scalar arithmetic
generator before taking that quotient.  Let
\[
 \mathcal J:\mathcal O(U)\longrightarrow\mathcal O(U)^2,
 \qquad \mathcal J f=(f,0),
\]
on the punctured $s$-domain $U$.  Since $R^2=-2sI$ and the displayed
connection is diagonal in the $(q_0,q_1)$ frame,
\[
 D\mathcal Jf=\mathcal J(f'),\qquad
 \widehat S\,\mathcal Jf=\mathcal J(sf),\qquad
 \widehat B_t\mathcal Jf
   =\mathcal J\!\left(sf-\frac t2f'\right).
\]
Consequently
\begin{equation}\label{eq:retained-branch-scalar-intertwiner}
 \widehat B_t\mathcal J=\mathcal J B_t,
 \qquad B_t=M-\frac t2\partial_s,
\end{equation}
as an equality of maps on the stated analytic-germ domain.  This is an
intertwiner for the even section only; it neither identifies the odd sheet
with a scalar function nor supplies descent of $D$ or $\widehat B_t$ to
the frozen quotient.  In particular, the branch commutator above remains
visible on the full two-component module.

\paragraph{Verification and continuation.}
The accompanying exact checker replays the determinant, inverse chart,
entire collision fibre, discriminant, coefficient heat flow, escaping
trajectory, and the fibre-product change of variables. It uses rational
polynomial arithmetic, not numerical evidence or a Lean certification.
The conductor and differential-module arguments above are full written
proofs. The next arithmetic step carries the actual heat and material
generators through the now-retained two-component arithmetic module,
including the branch coupling and analytic-unit derivative. No independent generic zero scan
is prescribed by this section.
